\chapter{Conclusions and Recommendations}\label{ch:concl}

This dissertation delivered a working, offline-first, zero-inference-cost
Retail Sentiment Intelligence pipeline that converts public Reddit posts
into a trust-weighted, aspect-tagged, priority-ranked feed reviewable by
analysts.  Against the four research questions posed in
Chapter~\ref{ch:problem}:

\begin{description}
\item[RQ1] A fine-tuned ModernBERT with a three-stage curriculum raised
      macro-F1 from 0.6272 to \textbf{0.7642} overall on the labelled sample
      and recovered the two long posts ($\ge 512$~tokens, $n{=}7$, all
      negative-class) the RoBERTa baseline mis-classifies (5/7
      $\to$~\textbf{7/7 correct}).
\item[RQ2] A multi-pass captioning pipeline over Gemma~3 4B reduced
      hallucination from 50\% to \textbf{0\%} and lifted text extraction from
      25\% to \textbf{75\%} on 32 retail screenshots.
\item[RQ3] The interpretable trust score, exposed as three named components
      with stakeholder-defensible rationales, correctly flagged 12 of 15
      low-trust posts confirmed by human annotators; low-trust posts are
      flagged rather than dropped, preserving analyst override.
\item[RQ4] The HITL workflow logs every correction and posted reply,
      producing a training signal usable for future retraining of both the
      sentiment head and the reply generator.
\end{description}

\section{Recommendations}
\begin{enumerate}
\item Adopt an incremental retraining cadence (monthly) using the HITL
      learning-loop store as the labelling stream.
\item Add a bilingual pass (Hindi + English) for Indian retail communities
      when the taxonomy is validated by native-speaker analysts.
\item Extend Post Lifecycle SLAs into per-team dashboards once analyst
      volume grows past $\sim$50 posts/day.
\end{enumerate}
